\chapter{\texorpdfstring{CONCLUSION}{}} %upper case only

This study has foregrounded the importance of storytelling as a critical tool in the creation and negotiation of identity, especially in times of extreme stress and uncertainty like immigration. Storytelling is one of the primary communication channels that molds and influences each layer of identity. As shown by earlier work that identities are created, socialized, and sustained, particularly in moments of uncertainty, transition, and coping through the practice of storytelling (Koenig Kellas 2022). This study concurs with and extends that understanding by foregrounding retrospective storytelling as a central mechanism in this process. The findings reveal that retrospective storytelling is not a deliberate or structured practice within families, but rather an embedded, everyday communicative activity through which meanings are continuously constructed. Both the content and the process of these narratives guide individual, relational, and intergenerational meaning-making of values and beliefs.

Importantly, the study also highlights the heightened influence of communal retrospective storytelling in shaping identity. In many instances, communal narratives circulating within close-knit Pakistani American (desi) communities play a decisive role in informing parental perspectives and actions. Parents draw upon these shared stories of challenges and moral expectations to guide major life decisions for their children. For example, data showed that some parents chose to homeschool their daughters through high school, not solely based on personal experience but in response to communal stories about other families whose children were perceived to have been “derailed” during high school becoming involved in behaviors  misaligned with Pakistani Muslim cultural and moral expectations, often associated with Westernized social norms. 

Impactful stories do not simply describe experience; they become part of us. Over time, individuals come to live their stories, drawing upon them to understand who they are and how they relate to others. Through the retelling of significant life episodes, individuals make sense of their past, negotiate their present identities, and project possible futures. In this way, storytelling functions not only as a means of remembering, but as an ongoing process of identity construction, sense-making, and self-understanding.

The findings suggest that retrospective storytelling does have a strong intergenerational impact on identity creation, and its stronger among immigrant families. Pakistani parents, mostly mothers share their stories of difficult times and challenges in early times of their marital relationships and during challenging circumstances with their children. Among children, the eldest child shares the most burden to process and cope with the trauma and unhealed wounds of the past or a household conflict, which in my data were mostly daughters. 

The findings support and extend the argument advanced by \citet{shinCommunicationTheoryIdentity2017}, who contend that when examining ethnic or racial identity groups, there is often greater variation within groups than between groups \citep[p.~3]{shinCommunicationTheoryIdentity2017}. Consistent with this claim, the data revealed substantial heterogeneity within the Pakistani American community. Even among Muslim participants, experiences, beliefs, family practices, and identity negotiations varied widely, challenging any monolithic understanding of Pakistani or Muslim identity. These challenges essentialized representations of Pakistani and Muslim identities and reinforces the need for more nuanced, within-group analyses in communication and identity scholarship.

The findings further demonstrated how territory and religion at birth intersect to shape identity formation, influencing personal and enacted identities through material conditions, communal affiliations, and ascribed relational expectations. Participants’ identities were not formed solely through individual choice, but through historically situated, culturally embedded, and relationally imposed meanings tied to nation, faith, migration status, and family roles. Together, these intersections illustrated how identity is dynamically negotiated across personal, enacted, relational, communal, and material layers, reinforcing the complexity and diversity that exists within Pakistani American identity formations.

Farah stated in her interview that she had dealt with a tough transition reconciling the tensions she perceived across the generations with each one torn between the guilt of not being able to be here (in the US) or there (Pakistan) simultaneously. Pakistan by constructing small Pakistani sub communities while others are lost in being a part of a no man zone where they neither fit here nor back home. Even if they try to have their own unique presence in the world on their own terms, they lose support and have no guidance, not knowing if things will work out or not. Since most first-generation parents don’t have a point of reference to look at.

Retrospective storytelling serves a dual role as a coping mechanism and a practice of recreation of ancestral culture in a foreign land. Therefore, the process of storytelling serves as a recollection site for the first-generation immigrant individuals where they recollect and narratively assemble the fragments of themselves from the past, making sense of who they were, who they became and who they want their children to be.

Stories also revealed the stress and dissonance that caused long-term mental stress manifesting the irreversible physical harm in the first generation. The children become a part of passive collateral damage as they see their parents in stress, arguing or worrying that they try to overcompensate for what their parents had missed because of the unfair amount of stress and anxiety caused by uncertainty. Moreover, demotion to a lower social class and economic status than they had achieved in Pakistan widened the distance between the world parents lived in prior to immigration, especially the mothers who had to uproot their lives, families, and dreams to move across the globe to support their husbands’ careers. The uncertainty and fears associated with immigration policies, visa status, and the financial burden of travel caused added more stress. 

Parents keep a close eye on the social circle of their children to monitor their friendships and social circles. In some cases, parents would draw a hardline on whom the girls could befriend. Girls are not encouraged to be friends with boys. Within the Muslim and Pakistani Muslim communities, parents restrict the social and friend circles of their children, depending on desi community surrounding the families. The parents might discourage their children from playing with peers they perceive to have greater freedom than their parents. Children in school settings often form connections with other migrant peers based on shared physical appearance and cultural overlap, with South Asian children frequently becoming sources of mutual support and informal networks for one another. In contrast, migrant parents sometimes experience tensions rooted in religious difference, which can produce boundaries that limit children’s participation in non-Muslim religious celebrations at friends’ homes or school-based events.

Intergenerational trauma is sustained at the personal layer through internalized expectations about gender, responsibility, and emotional endurance; at the relational layer through patterned parental guidance, reminders, and role based caregiving; and at the communal layer through culturally sanctioned narratives that position men’s career decisions as determinative of family destiny while women are cast in supportive roles. Retrospective storytelling functions as a CNSM sensemaking practice through which families explain ongoing hardship as necessary, inherited, or inevitable thereby stabilizing meaning across generations while simultaneously delaying disruption of traumatic family patterns.

These identity negotiations are not abstract; they are communicated and sustained through storytelling practices within families. Retrospective family stories function as key sites where intra-group differences, relational expectations, and historical inequalities are narrated, justified, and emotionally transferred across generations. Through repeated storytelling of hardship, migration, and unresolved conflict, trauma becomes embedded in family narratives and incorporated into children’s sensemaking processes. In this way, storytelling operates as a central mechanism through which identity diversity, power asymmetries, and unhealed wounds are not only remembered but also transmitted, negotiated, and, at times, reenacted across generations.

By centering Pakistani migrant families, this study extends CNSM and CTI into non-White, transnational contexts, demonstrating how retrospective storytelling operates as a culturally situated sense-making practice that simultaneously sustains communal coherence and reproduces gendered and intergenerational power dynamics across personal, relational, and communal identity layers.

\section{Future Directions for Research}

While the present study offers an initial exploration into how Pakistani American immigrants negotiate their identity construction and navigate the shifting dynamics of intergenerational roles, it represents only a partial mapping of a complex sociocultural landscape. 

There are myriad opportunities for further understanding the complex nature of the process which follows the patterns of the previous generation and history of the country. The future scholarship should expand upon several critical avenues identified in this research.

\subsection{Socio-Economic Stratification and its presentation in the Material Frame}

A running theme throughout this study was exploration of the role of the material frame in shaping the experience of first, one-and-a-half and second-generation Pakistani Americans. It became clear that this experience is guided by the socio-economic status the first-generation immigrants experienced in Pakistan. This, then, highlights a potential status knowledge gap between the access immigrants enjoyed in Pakistan versus their initial experiences in the US Future research should utilize a longitudinal approach to examine how the socio-economic status (SES) of origin interacts with the SES attained in the host country.

\subsection{The ``Zeroth Generation''}

To gain a truly holistic view of the immigrant trajectory, researchers should conduct interviews with the “zeroth generation” (the parents of first-generation immigrants). Understanding their perspectives would clarify how ancestral expectations and back-home comparisons exert pressure on the identity formation of their children and grandchildren.

\subsection{Intra-Ethnic Hierarchies and Professional Networks}

This study underscores that the Pakistani expatriate community is not a monolith but is subdivided by a series of ethno-professional socio-economic strata. Similar to the “Status Knowledge Gap,” the socio-economic stratum each immigrant joined, based on their regional, linguistic, ethnic, educational, and professional credentials, also played a significant role in the advice, experiences and communal stories they were exposed to and internalized. Future studies should investigate how these filtered networks can impact the immigrant experience.

\subsection{Comparative Racialization and Visibility}

To better understand the specificities of the Pakistani American experience, subsequent research should include comparative analyses involving mixed-race couples and white (European) immigrants. This would allow for a more nuanced understanding of how visible differences and racialized materialities influence social integration.

\section{Practical Implications}

The findings in this study provide us with a multitude of directions to utilize the translational heuristics of CNSM and develop materials to support this underserved population.

There is a demonstrated need for guidance documents or resource books designed for families navigating this sea of role-confusions, role-reversals, and identity gaps. These tools could help immigrants who feel lost in the transition to better support their children’s educational and social development.

Parental involvement should be recognized as a valuable resource in cultivating culturally responsive classrooms. Schools can facilitate this by providing structured opportunities for parents to participate in classroom activities, contribute educational materials, or co-organize events that reflect diverse cultural perspectives. Supporting parents in sharing brief, accessible, and interactive content such as informational booklets or age-appropriate activities can enhance engagement while fostering cross-cultural understanding.On the other end of the spectrum, educational institutions and social service providers can use these insights to inform policies that recognize the multilayered identities of South Asian students. By understanding the pressures of inter-community comparison and the lack of ethnic social capital (e.g., the absence of desi peer groups in certain districts), practitioners can better facilitate acculturation.

The findings underscore the need for institutional policies that intentionally support culturally inclusive and responsive educational environments. At the structural level, school districts should prioritize the recruitment and retention of a diverse workforce, including teachers, administrators, and support staff who bring varied cultural perspectives and demonstrate cultural competence. In addition, ongoing professional development should be mandated to equip educators with the skills necessary to engage effectively with culturally and linguistically diverse students and families. Such training should move beyond surface-level awareness to include reflexive practices, anti-bias education, and strategies for inclusive communication.

At the student level, these inclusive practices contribute to safer learning environments where children feel validated in expressing their identities. When students see their cultural backgrounds reflected and respected within institutional spaces, they are more likely to develop confidence in their sense of self and experience reduced identity tensions. Thus, schools function not only as sites of academic learning but also as critical spaces for identity affirmation and intercultural dialogue.

The educational institutions should adopt a holistic approach that recognizes schools as key sites of identity development. Policies that support inclusive environments, culturally responsive teaching, and active family engagement can reduce identity tensions and foster a stronger sense of belonging among students from immigrant backgrounds. In doing so, schools can play a transformative role in supporting equitable educational experiences and more integrated identity development across diverse populations.

\section{Limitations}

It is important to acknowledge that this study is far from a comprehensive and fully representative model for the Pakistani American immigrant population. There are limitations to what was able to be accomplished by this study and the length to which any conclusions drawn from this study can be taken.

Most importantly, this is a qualitative and narrative study. While this allows for deep exploration of individual experiences, it does not lend itself to a generalizable model of immigrant experience. Thus, the poems and the conclusions of this study should be taken as individual, microcosmic representations of individuals, rather than a generalizable lens with which you can cover the full breadth of Pakistani American immigrants.

Additionally, given the diversity of population in Pakistan, it was nigh impossible to cover a cross section of the entire population. A significant limitation in this study was that regretfully, I was unable to include any non-Muslim Pakistani American immigrants. Their inclusion would have provided an additional perspective given the importance of religion in Pakistani American identity.

\section{Personal Narrative: Living the Afterlife of Immigrants}

With the suspension of immigrant visas for 75 countries, including Pakistan, the political landscape has reactivated historical patterns of exclusion that many of us believed belonged to another era. What I have spent years studying; immigrant lived experiences, parenting under precarity, and the emotional labor of raising children amid political uncertainty is no longer confined to interviews or analysis. It is unfolding in real time, within our surroundings.

We chose migration, believing in education, mobility, and the promise of opportunity. We left home voluntarily, trusting that our movement was forward, not cyclical. Now, my children, with teary eyes and emotions far larger than their years, are preparing to leave the only home they have known for the last seven and five years. What was meant to be stability has once again become shaky ground.

The most painful moments are not procedural; they are intimate. I have seen fear and confusion in my seven-year-old’s eyes as she asks why they do not want us here. That question carries generations of displacement. It reminds me that immigration policy is never abstract. It enters classrooms, homes, and childhoods, shaping how children come to understand belonging, safety, and worth.

Fear-mongering is no longer distant or imagined. Stories of children taken from schools, of people detained from roads based on accents and appearances, and of people detained while they were attending their immigration court hearings circulate widely. These stories attach fear to bodies and turn difference into danger. Another political layer of trauma is quietly becoming part of immigrant identity, replacing the hope of becoming global citizens with the reality of conditional belonging and constant vigilance.
This research began as an effort to understand immigrant lived experiences and parenting across generations. It concludes with the realization that (im)migration is not a single journey, but a recurring condition. Borders change, policies shift, but the emotional consequences endure. The stories shared by participants, alongside my own, reveal how identity is continually shaped by rupture, memory, and forced movement rather than arrival.

To conclude, this work is not to resolve these tensions, but to name them. Immigration is lived through history, carried into the present, and inherited by children who did not choose it. Scholarship, like migration, does not stand outside life. It is entangled with it written not only in data and analysis, but in fear, loss, love, and the quiet work of holding children together when the ground beneath them shifts again.
